\section{Syzygies, reversal, and where a second brake can occur}\label{sec:reversal}

Time reversal about a second brake, combined with Montgomery's theorems on
syzygies, restricts where a second brake can lie. A syzygy is a collinear
instant. Every syzygy of a zero-angular-momentum orbit launched from a
noncollinear rest configuration is transverse, and every brake of such an
orbit is noncollinear (Section~\ref{sec:search}); so syzygies are isolated.

\begin{theorem}[Montgomery~\cite{Montgomery2002}]
Every bounded, collision-free, zero-angular-momentum solution of the planar
three-body problem, with any positive masses, has infinitely many syzygies.
Moreover the normalized signed area
\[
 z=\frac{4\sqrt3\,\Delta}{r_{12}^2+r_{23}^2+r_{31}^2},
 \qquad \Delta=\tfrac12(q_2-q_1)\times(q_3-q_1),
\]
satisfies $\frac{d}{dt}(f\dot z)=-qz$ along every zero-angular-momentum
solution away from triple collision, with $f>0$ and $q\ge0$, where $q$
vanishes only on the Lagrange homothetic solutions.
\end{theorem}

Suppose the tied orbit brakes again at $\tau>0$ after a collision-free
segment. Uniqueness and reversal give $q(\tau+s)=q(\tau-s)$ and
$q(2\tau+s)=q(s)$: the orbit is $2\tau$-periodic, bounded, and
collision-free for all time. Let $0<s_1<\dots<s_n<\tau$ be its syzygies
before $\tau$. Reflection sends them to the syzygies $2\tau-s_k$ of
$(\tau,2\tau)$, with $q(2\tau-s_k)=q(s_k)$ and
$\dot q(2\tau-s_k)=-\dot q(s_k)$.

\begin{theorem}[No brake before the first syzygy]\label{thm:first-syzygy}
For every real $u$, and for every noncollinear brake initial condition with
any positive masses, there is no second brake at or before the first
positive syzygy time.
\end{theorem}
\begin{proof}
A brake is noncollinear, so a brake at or before the first syzygy occurs
strictly before it, with $n=0$. By reflection the periodic orbit has no
syzygy in $(\tau,2\tau)$ either, nor at $0$, $\tau$, $2\tau$. It is bounded,
collision-free, and has zero angular momentum, yet has no syzygy at all,
contradicting Montgomery's theorem.
\end{proof}

This settles unconditionally the first-arc objective of the torque-history
route of Section~\ref{sec:torque}; the corrected conditional lemma and its
terminal hypothesis are not needed for that purpose.

\begin{theorem}[A second brake is a stutter]\label{thm:stutter}
If the orbit brakes at $\tau$ with $n\ge1$ earlier syzygies, the two
syzygies adjacent to $\tau$, at $s_n$ and $2\tau-s_n$, have the same
collinear configuration and opposite velocities; in particular the same
middle body. The syzygy sequence of the period is the double palindrome
$(\sigma_1\cdots\sigma_n)(\sigma_n\cdots\sigma_1)$ repeated, so the forward
sequence contains a repeated adjacent letter (a \emph{stutter}, in the
terminology of~\cite{MMV}) at every brake.
\end{theorem}

Hence an inter-syzygy arc whose two syzygies have different middle bodies
contains no brake, without any further estimate. This is the converse of
the observation in~\cite{MMV} that brake orbits run backwards produce an
open set of stuttering orbits.

\begin{theorem}[Position within the arc]\label{thm:zmax}
Between consecutive syzygies $z$ has exactly one critical point, a strict
extremum. A second brake lies at that extremum of its stutter arc, and $I$
has a strict maximum there.
\end{theorem}
\begin{proof}
Where $z>0$, $f\dot z$ is nonincreasing, so $\dot z$ changes sign at most
once, and $\dot z\equiv0$ on an interval would force $qz=0$, that is a
Lagrange homothetic motion. All velocities vanish at a brake, so
$\dot z(\tau)=0$; the inertia statement is Lagrange--Jacobi.
\end{proof}

\begin{theorem}[Brakes as fixed points on the syzygy section]\label{thm:involution}
Let $\Sigma$ be the transverse syzygy states, $P$ the next-syzygy map on its
domain in $\Sigma$, and $R(q,v)=(q,-v)$. Then $RPR=P^{-1}$, so $RP$ is an
involution, and the inter-syzygy arc starting at $\sigma\in\Sigma$ contains a
brake if and only if $RP(\sigma)=\sigma$.
\end{theorem}
\begin{proof}
Reversibility gives $RPR=P^{-1}$. If $P(\sigma)=R\sigma$ on an arc of
duration $T$, then $q(T)=q(0)$ and $v(T)=-v(0)$; uniqueness applied to
$t\mapsto q(T-t)$ gives $q(T/2+s)=q(T/2-s)$, hence $v(T/2)=0$. The converse
is Theorem~\ref{thm:stutter}.
\end{proof}

The fixed set of $RP$ is exactly the set of syzygy states of brake-to-syzygy
arcs run backwards, that is, the image of the syzygy map of~\cite{MMV}
together with its velocities. For the tied family the second-brake question
becomes an intersection problem between the one-parameter curves
$u\mapsto\sigma_k(u)$ of $k$th syzygy states and this two-dimensional fixed
set inside the four-dimensional reduced section. The codimension count is
unchanged, but the brake-reachable states are now an explicit geometric
object, and Conjecture~1 of~\cite{MMV} (that $\dot I<0$ from a brake through
its first syzygy) would add that $I$ is strictly monotone on each half of a
brake's stutter arc.

These statements are consistent with the global numerical picture of the
tied family. In a scan of about $170{,}000$ parameters on $[0.10,0.4142]$
with steps down to $10^{-6}$, every inter-syzygy arc had exactly one
$z$-extremum, the first inertia minimum never preceded the first syzygy,
about half of the escaping trajectories had no stutter before the escape
lemma of Section~\ref{sec:events} applied, and all of the forty closest
approaches of the family to a second brake occurred on stutter arcs, as
Theorem~\ref{thm:stutter} predicts by continuity. The closest approaches found,
after refinement to about $10^{-13}$ in $u$, are at $u\simeq0.2003099995$,
$t\simeq7.91$, and $u\simeq0.2556321826$, $t\simeq4.79$, where the kinetic
energy at its local minimum is only $7\times10^{-9}\,U_0$ and
$8\times10^{-9}\,U_0$ respectively (residual velocities of order
$2\times10^{-4}$); neither is a zero, and neither lies near a rational of
small denominator beyond chance level. Section~\ref{sec:search}
uses these near misses as seeds.
